\documentclass{amsart}
% For dealing with references we use the comment environment
\usepackage{verbatim}
\newenvironment{reference}{\comment}{\endcomment}
%\newenvironment{reference}{}{}
\newenvironment{slogan}{\comment}{\endcomment}
\newenvironment{history}{\comment}{\endcomment}

% For commutative diagrams we use Xy-pic
\usepackage[all]{xy}

% We use 2cell for 2-commutative diagrams.
\xyoption{2cell}
\UseAllTwocells

% We use multicol for the list of chapters between chapters
\usepackage{multicol}

% This is generally recommended for better output
\usepackage{lmodern}
\usepackage[T1]{fontenc}

% For cross-file-references

% Package for hypertext links:
\usepackage{hyperref}

% For any local file, say "hello.tex" you want to link to please

% Theorem environments.
%
\theoremstyle{plain}
\newtheorem{theorem}[subsection]{Theorem}
\newtheorem{proposition}[subsection]{Proposition}
\newtheorem{lemma}[subsection]{Lemma}

\theoremstyle{definition}
\newtheorem{definition}[subsection]{Definition}
\newtheorem{example}[subsection]{Example}
\newtheorem{exercise}[subsection]{Exercise}
\newtheorem{situation}[subsection]{Situation}

\theoremstyle{remark}
\newtheorem{remark}[subsection]{Remark}
\newtheorem{remarks}[subsection]{Remarks}

\numberwithin{equation}{subsection}

% Macros
%
\def\lim{\mathop{\mathrm{lim}}\nolimits}
\def\colim{\mathop{\mathrm{colim}}\nolimits}
\def\Spec{\mathop{\mathrm{Spec}}}
\def\Hom{\mathop{\mathrm{Hom}}\nolimits}
\def\Ext{\mathop{\mathrm{Ext}}\nolimits}
\def\SheafHom{\mathop{\mathcal{H}\!\mathit{om}}\nolimits}
\def\SheafExt{\mathop{\mathcal{E}\!\mathit{xt}}\nolimits}
\def\Sch{\mathit{Sch}}
\def\Mor{\mathop{\mathrm{Mor}}\nolimits}
\def\Ob{\mathop{\mathrm{Ob}}\nolimits}
\def\Sh{\mathop{\mathit{Sh}}\nolimits}
\def\NL{\mathop{N\!L}\nolimits}
\def\CH{\mathop{\mathrm{CH}}\nolimits}
\def\proetale{{pro\text{-}\acute{e}tale}}
\def\etale{{\acute{e}tale}}
\def\QCoh{\mathit{QCoh}}
\def\Ker{\mathop{\mathrm{Ker}}}
\def\Im{\mathop{\mathrm{Im}}}
\def\Coker{\mathop{\mathrm{Coker}}}
\def\Coim{\mathop{\mathrm{Coim}}}

% Boxtimes
%
\DeclareMathSymbol{\boxtimes}{\mathbin}{AMSa}{"02}

%
% Macros for moduli stacks/spaces
%
\def\QCohstack{\mathcal{QC}\!\mathit{oh}}
\def\Cohstack{\mathcal{C}\!\mathit{oh}}
\def\Spacesstack{\mathcal{S}\!\mathit{paces}}
\def\Quotfunctor{\mathrm{Quot}}
\def\Hilbfunctor{\mathrm{Hilb}}
\def\Curvesstack{\mathcal{C}\!\mathit{urves}}
\def\Polarizedstack{\mathcal{P}\!\mathit{olarized}}
\def\Complexesstack{\mathcal{C}\!\mathit{omplexes}}
% \Pic is the operator that assigns to X its picard group, usage \Pic(X)
% \Picardstack_{X/B} denotes the Picard stack of X over B
% \Picardfunctor_{X/B} denotes the Picard functor of X over B
\def\Pic{\mathop{\mathrm{Pic}}\nolimits}
\def\Picardstack{\mathcal{P}\!\mathit{ic}}
\def\Picardfunctor{\mathrm{Pic}}
\def\Deformationcategory{\mathcal{D}\!\mathit{ef}}

\setlength{\emergencystretch}{3em}
\begin{document}
\title[Stacks Project, Tag 09S8]{1323 --- Rickard tilting-object criterion for derived equivalence between an ordinary algebra and a differential graded algebra}
\maketitle
\noindent Copyright (C) 2005--2025 Johan de Jong. Stacks Project authors. Source: \href{https://stacks.math.columbia.edu/tag/09S8}{Tag 09S8}.

\noindent Editorial difficulty note (not author text). Difficulty H2: The author converts an abstract compact generating object with concentrated self-Ext into an actual differential graded endomorphism algebra, truncates that algebra through a quasi-isomorphism zigzag, and obtains the derived equivalence carrying the ordinary algebra to the chosen object. This enhancement and tilting construction is substantially beyond ordinary qualification algebra..

\noindent Editorial provenance note (not author text). History: excerpt from revision \texttt{a0444\allowbreak 6e57e\allowbreak c1fbc\allowbreak 252a8\allowbreak 71afc\allowbreak ec775\allowbreak 2fb28\allowbreak 07b14}, dga.tex, lines 6584--6645. Reference presentation was adapted; original-author context and core support blocks were retained or added. The mathematical wording is unchanged. Licensed under GNU FDL 1.2 or later; no invariant sections or cover texts. See ../licenses/COPYING. Context blocks reproduce author definitions and supporting results; they are not additional counted problems.

\begin{definition}
\label{definition-dga}
Let $R$ be a commutative ring. A {\it differential graded algebra over $R$}
is either
\begin{enumerate}
\item a chain complex $A_\bullet$ of $R$-modules endowed with
$R$-bilinear maps $A_n \times A_m \to A_{n + m}$,
$(a, b) \mapsto ab$ such that
$$
\text{d}_{n + m}(ab) = \text{d}_n(a)b + (-1)^n a\text{d}_m(b)
$$
and such that $\bigoplus A_n$ becomes an associative and unital
$R$-algebra, or
\item a cochain complex $A^\bullet$ of $R$-modules endowed with
$R$-bilinear maps $A^n \times A^m \to A^{n + m}$, $(a, b) \mapsto ab$
such that
$$
\text{d}^{n + m}(ab) = \text{d}^n(a)b + (-1)^n a\text{d}^m(b)
$$
and such that $\bigoplus A^n$ becomes an associative and unital $R$-algebra.
\end{enumerate}
\end{definition}

\begin{definition}
\label{definition-dgm}
Let $R$ be a ring.
Let $(A, \text{d})$ be a differential graded algebra over $R$.
A (right) {\it differential graded module} $M$ over $A$ is a right $A$-module
$M$ which has a grading $M = \bigoplus M^n$ and a differential $\text{d}$
such that $M^n A^m \subset M^{n + m}$, such that
$\text{d}(M^n) \subset M^{n + 1}$, and such that
$$
\text{d}(ma) = \text{d}(m)a + (-1)^n m\text{d}(a)
$$
for $a \in A$ and $m \in M^n$. A
{\it homomorphism of differential graded modules} $f : M \to N$
is an $A$-module map compatible with gradings and differentials.
The category of (right) differential graded $A$-modules is denoted
$\text{Mod}_{(A, \text{d})}$.
\end{definition}

\begin{lemma}
\label{lemma-resolve}
Let $(A, \text{d})$ be a differential graded algebra.
Let $M$ be a differential graded $A$-module. There exists a homomorphism
$P \to M$ of differential graded $A$-modules such that
\begin{enumerate}
\item $P \to M$ is a quasi-isomorphism, and
\item $P$ has property (P).
\end{enumerate}
\end{lemma}

\begin{proof}
Set $M = M_0$. We inductively choose short exact sequences
$$
0 \to M_{i + 1} \to P_i \to M_i \to 0
$$
where the maps $P_i \to M_i$ are chosen as in Lemma \ref{lemma-good-quotient}.
This gives a ``resolution''
$$
\ldots \to P_2 \xrightarrow{f_2} P_1 \xrightarrow{f_1} P_0 \to M \to 0
$$
Then we set
$$
P = \bigoplus\nolimits_{i \geq 0} P_i
$$
as an $A$-module with grading given by
$P^n = \bigoplus_{a + b = n} P_{-a}^b$ and
differential (as in the construction of the total complex associated
to a double complex) by
$$
\text{d}_P(x) = f_{-a}(x) + (-1)^a \text{d}_{P_{-a}}(x)
$$
for $x \in P_{-a}^b$. With these conventions $P$ is indeed a differential
graded $A$-module. Recalling that each $P_i$ has a two step filtration
$0 \to P_i' \to P_i \to P_i'' \to 0$ we set
$$
F_{2i}P = \bigoplus\nolimits_{i \geq j \geq 0} P_j
\subset
\bigoplus\nolimits_{i \geq 0} P_i = P
$$
and we add $P'_{i + 1}$ to $F_{2i}P$ to get $F_{2i + 1}$.
These are differential graded submodules and the successive quotients
are direct sums of shifts of $A$. By
Lemma \href{https://stacks.math.columbia.edu/tag/09K0}{lemma target graded projective (Tag 09K0)} we see that
the inclusions $F_iP \to F_{i + 1}P$ are admissible monomorphisms.
Finally, we have to show that the map $P \to M$ (given by the
augmentation $P_0 \to M$) is a quasi-isomorphism. This follows from
Homology, Lemma \href{https://stacks.math.columbia.edu/tag/09IZ}{lemma good resolution gives qis (Tag 09IZ)}.
\end{proof}


\begin{lemma}
\label{lemma-good-quotient}
Let $(A, \text{d})$ be a differential graded algebra.
Let $M$ be a differential graded $A$-module. There exists a homomorphism
$P \to M$ of differential graded $A$-modules with the following
properties
\begin{enumerate}
\item $P \to M$ is surjective,
\item $\Ker(\text{d}_P) \to \Ker(\text{d}_M)$ is surjective, and
\item $P$ sits in an admissible short exact sequence
$0 \to P' \to P \to P'' \to 0$ where $P'$, $P''$ are direct sums
of shifts of $A$.
\end{enumerate}
\end{lemma}

\begin{proof}
Let $P_k$ be the free $A$-module with generators $x, y$ in degrees
$k$ and $k + 1$. Define the structure of a differential graded
$A$-module on $P_k$ by setting $\text{d}(x) = y$ and $\text{d}(y) = 0$.
For every element $m \in M^k$ there is a homomorphism
$P_k \to M$ sending $x$ to $m$ and $y$ to $\text{d}(m)$.
Thus we see that there is a surjection from a direct sum
of copies of $P_k$ to $M$. This clearly produces $P \to M$
having properties (1) and (3). To obtain property (2) note
that if $m \in \Ker(\text{d}_M)$ has degree $k$, then there is a map
$A[k] \to M$ mapping $1$ to $m$. Hence we can achieve (2) by adding
a direct sum of copies of shifts of $A$.
\end{proof}


\begin{lemma}
\label{lemma-property-P-sequence}
Let $(A, \text{d})$ be a differential graded algebra.
Let $P$ be a differential graded $A$-module. If $F_\bullet$
is a filtration as in property (P), then we obtain an
admissible short exact sequence
$$
0 \to
\bigoplus\nolimits F_iP \to
\bigoplus\nolimits F_iP \to P \to 0
$$
of differential graded $A$-modules.
\end{lemma}

\begin{proof}
The second map is the direct sum of the inclusion maps.
The first map on the summand $F_iP$ of the source is the sum
of the identity $F_iP \to F_iP$ and the negative of the inclusion
map $F_iP \to F_{i + 1}P$. Choose homomorphisms $s_i : F_{i + 1}P \to F_iP$
of graded $A$-modules which are left inverse to the inclusion maps.
Composing gives maps $s_{j, i} : F_jP \to F_iP$ for all $j > i$.
Then a left inverse of the first arrow maps $x \in F_jP$ to
$(s_{j, 0}(x), s_{j, 1}(x), \ldots, s_{j, j - 1}(x), 0, \ldots)$
in $\bigoplus F_iP$.
\end{proof}


\begin{lemma}
\label{lemma-property-P-K-projective}
Let $(A, \text{d})$ be a differential graded algebra.
Let $P$ be a differential graded $A$-module with property (P).
Then
$$
\Hom_{K(\text{Mod}_{(A, \text{d})})}(P, N) = 0
$$
for all acyclic differential graded $A$-modules $N$.
\end{lemma}

\begin{proof}
We will use that $K(\text{Mod}_{(A, \text{d})})$ is a triangulated
category (Proposition \href{https://stacks.math.columbia.edu/tag/09KJ}{proposition homotopy category triangulated (Tag 09KJ)}).
Let $F_\bullet$ be a filtration on $P$ as in property (P).
The short exact sequence of Lemma \ref{lemma-property-P-sequence}
produces a distinguished triangle. Hence by
Derived Categories, Lemma \href{https://stacks.math.columbia.edu/tag/0149}{lemma representable homological (Tag 0149)}
it suffices to show that
$$
\Hom_{K(\text{Mod}_{(A, \text{d})})}(F_iP, N) = 0
$$
for all acyclic differential graded $A$-modules $N$ and all $i$.
Each of the differential graded modules $F_iP$ has a finite filtration
by admissible monomorphisms, whose graded pieces are direct sums
of shifts $A[k]$. Thus it suffices to prove that
$$
\Hom_{K(\text{Mod}_{(A, \text{d})})}(A[k], N) = 0
$$
for all acyclic differential graded $A$-modules $N$ and all $k$.
This follows from Lemma \href{https://stacks.math.columbia.edu/tag/09K1}{lemma hom from shift free (Tag 09K1)}.
\end{proof}


\begin{lemma}
\label{lemma-hom-derived}
Let $(A, \text{d})$ be a differential graded algebra.
Let $M$ and $N$ be differential graded $A$-modules.
\begin{enumerate}
\item Let $P \to M$ be a P-resolution as in
Lemma \ref{lemma-resolve}. Then
$$
\Hom_{D(A, \text{d})}(M, N) =
\Hom_{K(\text{Mod}_{(A, \text{d})})}(P, N)
$$
\item Let $N \to I$ be an I-resolution as in
Lemma \href{https://stacks.math.columbia.edu/tag/09KU}{lemma right resolution (Tag 09KU)}. Then
$$
\Hom_{D(A, \text{d})}(M, N) =
\Hom_{K(\text{Mod}_{(A, \text{d})})}(M, I)
$$
\end{enumerate}
\end{lemma}

\begin{proof}
Let $P \to M$ be as in (1). Since $P \to M$ is a quasi-isomorphism we see that
$$
\Hom_{D(A, \text{d})}(P, N) = \Hom_{D(A, \text{d})}(M, N)
$$
by definition of the derived category. A morphism
$f : P \to N$ in $D(A, \text{d})$ is equal to
$s^{-1}f'$ where $f' : P \to N'$ is a morphism and
$s : N \to N'$ is a quasi-isomorphism. Choose a distinguished triangle
$$
N \to N' \to Q \to N[1]
$$
As $s$ is a quasi-isomorphism, we see that $Q$ is acyclic. Thus
$\Hom_{K(\text{Mod}_{(A, \text{d})})}(P, Q[k]) = 0$ for all $k$ by
Lemma \ref{lemma-property-P-K-projective}. Since
$\Hom_{K(\text{Mod}_{(A, \text{d})})}(P, -)$
is cohomological, we conclude that we can lift $f' : P \to N'$
uniquely to a morphism $f : P \to N$. This finishes the proof.

\medskip\noindent
The proof of (2) is dual to that of (1) using
Lemma \href{https://stacks.math.columbia.edu/tag/09KS}{lemma property I K injective (Tag 09KS)} in stead of
Lemma \ref{lemma-property-P-K-projective}.
\end{proof}


\begin{lemma}
\label{lemma-tilting-equivalence}
Let $R$ be a ring. Let $(A, \text{d})$ and $(B, \text{d})$ be 
differential graded algebras over $R$. Let $N$ be a
differential graded $(A, B)$-bimodule. Assume that
\begin{enumerate}
\item $N$ defines a compact object of $D(B, \text{d})$,
\item if $N' \in D(B, \text{d})$ and
$\Hom_{D(B, \text{d})}(N, N'[n]) = 0$ for $n \in \mathbf{Z}$,
then $N' = 0$, and
\item the map $H^k(A) \to \Hom_{D(B, \text{d})}(N, N[k])$ is an
isomorphism for all $k \in \mathbf{Z}$.
\end{enumerate}
Then
$$
- \otimes_A^\mathbf{L} N : D(A, \text{d}) \to D(B, \text{d})
$$
gives an $R$-linear equivalence of triangulated categories.
\end{lemma}

\begin{proof}
By Lemma \ref{lemma-tensor-with-compact-fully-faithful}
the functor $M \longmapsto M \otimes_A^\mathbf{L} N$ is
fully faithful. By Lemma \href{https://stacks.math.columbia.edu/tag/09LT}{lemma tensor hom adjoint (Tag 09LT)}
the functor $N' \longmapsto R\Hom(N, N')$ is a right adjoint.
By assumption (3) the kernel of $R\Hom(N, -)$ is zero.
Hence the result follows from
Derived Categories, Lemma
\href{https://stacks.math.columbia.edu/tag/09J1}{lemma fully faithful adjoint kernel zero (Tag 09J1)}.
\end{proof}


\begin{lemma}
\label{lemma-tensor-with-compact-fully-faithful}
With notation and assumptions as in Lemma \href{https://stacks.math.columbia.edu/tag/09LT}{lemma tensor hom adjoint (Tag 09LT)}.
Assume
\begin{enumerate}
\item $N$ defines a compact object of $D(B, \text{d})$, and
\item the map $H^k(A) \to \Hom_{D(B, \text{d})}(N, N[k])$ is an
isomorphism for all $k \in \mathbf{Z}$.
\end{enumerate}
Then the functor $-\otimes_A^\mathbf{L} N$ is fully faithful.
\end{lemma}

\begin{proof}
Our functor has a left adjoint given by
$R\Hom(N, -)$ by Lemma \href{https://stacks.math.columbia.edu/tag/09LT}{lemma tensor hom adjoint (Tag 09LT)}.
By Categories, Lemma \href{https://stacks.math.columbia.edu/tag/07RB}{lemma adjoint fully faithful (Tag 07RB)}
it suffices to show that for a differential graded $A$-module $M$
the map
$$
M \longrightarrow R\Hom(N, M \otimes_A^\mathbf{L} N)
$$
is an isomorphism in $D(A, \text{d})$. For this it suffices to show that
$$
H^n(M) \longrightarrow
\text{Ext}^n_{D(B, \text{d})}(N, M \otimes_A^\mathbf{L} N)
$$
is an isomorphism, see Lemma \href{https://stacks.math.columbia.edu/tag/0CS5}{lemma derived restriction exts (Tag 0CS5)}.
Since $N$ is a compact object the right hand side commutes
with direct sums. Thus by Remark \href{https://stacks.math.columbia.edu/tag/0FQE}{remark P resolution (Tag 0FQE)}
it suffices to prove this map is an isomorphism for $M = A[k]$.
Since $(A[k] \otimes_A^\mathbf{L} N) = N[k]$ by
Remark \href{https://stacks.math.columbia.edu/tag/0FQN}{remark shift tensor no sign (Tag 0FQN)},
assumption (2) on $N$ is that the result holds for these.
\end{proof}


\begin{lemma}
\label{lemma-qis-equivalence}
Let $R$ be a ring. Let $(A, \text{d}) \to (B, \text{d})$ be a
homomorphism of differential graded algebras over $R$, which induces
an isomorphism on cohomology algebras. Then
$$
- \otimes_A^\mathbf{L} B : D(A, \text{d}) \to D(B, \text{d})
$$
gives an $R$-linear equivalence of triangulated categories with
quasi-inverse the restriction functor $N \mapsto N_A$.
\end{lemma}

\begin{proof}
By Lemma \ref{lemma-tensor-with-compact-fully-faithful}
the functor $M \longmapsto M \otimes_A^\mathbf{L} B$ is
fully faithful. By Lemma \href{https://stacks.math.columbia.edu/tag/09LT}{lemma tensor hom adjoint (Tag 09LT)}
the functor $N \longmapsto R\Hom(B, N) = N_A$ is a right adjoint, see
Example \href{https://stacks.math.columbia.edu/tag/0BYX}{example map hom tensor (Tag 0BYX)}.
It is clear that the kernel of $R\Hom(B, -)$ is zero.
Hence the result follows from
Derived Categories, Lemma
\href{https://stacks.math.columbia.edu/tag/09J1}{lemma fully faithful adjoint kernel zero (Tag 09J1)}.
\end{proof}


\section{P-resolutions}
\label{section-P-resolutions}

\noindent
This section is the analogue of
Derived Categories, Section \href{https://stacks.math.columbia.edu/tag/06XW}{section unbounded (Tag 06XW)}.

\medskip\noindent
Let $(A, \text{d})$ be a differential graded algebra.
Let $P$ be a differential graded $A$-module. We say $P$
{\it has property (P)} if it there exists a filtration
$$
0 = F_{-1}P \subset F_0P \subset F_1P \subset \ldots \subset P
$$
by differential graded submodules such that
\begin{enumerate}
\item $P = \bigcup F_pP$,
\item the inclusions $F_iP \to F_{i + 1}P$ are admissible
monomorphisms,
\item the quotients $F_{i + 1}P/F_iP$ are isomorphic as differential
graded $A$-modules to a direct sum of $A[k]$.
\end{enumerate}
In fact, condition (2) is a consequence of condition (3), see
Lemma \href{https://stacks.math.columbia.edu/tag/09K0}{lemma target graded projective (Tag 09K0)}. Moreover, the reader
can verify that as a graded $A$-module $P$ will be isomorphic to a
direct sum of shifts of $A$.


\begin{lemma}
\label{lemma-rickard}
Let $R$ be a ring. Let $(A, \text{d})$ and $(B, \text{d})$ be
differential graded $R$-algebras. Assume that $A = H^0(A)$.
The following are equivalent
\begin{enumerate}
\item $D(A, \text{d})$ and $D(B, \text{d})$ are equivalent as $R$-linear
triangulated categories, and
\item there exists an object $P$ of $D(B, \text{d})$ such that
\begin{enumerate}
\item $P$ is a compact object of $D(B, \text{d})$,
\item if $N \in D(B, \text{d})$ with $\Hom_{D(B, \text{d})}(P, N[i]) = 0$
for $i \in \mathbf{Z}$, then $N = 0$,
\item $\Hom_{D(B, \text{d})}(P, P[i]) = 0$ for $i \not = 0$ and
equal to $A$ for $i = 0$.
\end{enumerate}
\end{enumerate}
The equivalence $D(A, \text{d}) \to D(B, \text{d})$
constructed in (2) sends $A$ to $P$.
\end{lemma}

\begin{proof}
Let $F : D(A, \text{d}) \to D(B, \text{d})$ be an equivalence.
Then $F$ maps compact objects to compact objects. Hence $P = F(A)$ is
compact, i.e., (2)(a) holds. Conditions (2)(b) and (2)(c) are immediate
from the fact that $F$ is an equivalence.

\medskip\noindent
Let $P$ be an object as in (2). Represent $P$ by a
differential graded module with property (P). Set
$$
(E, \text{d}) = \Hom_{\text{Mod}^{dg}_{(B, \text{d})}}(P, P)
$$
Then $H^0(E) = A$ and $H^k(E) = 0$ for $k \not = 0$ by
Lemma \ref{lemma-hom-derived} and assumption (2)(c).
Viewing $P$ as a $(E, B)$-bimodule and using
Lemma \ref{lemma-tilting-equivalence} and assumption (2)(b)
we obtain an equivalence
$$
D(E, \text{d}) \to D(B, \text{d})
$$
sending $E$ to $P$.
Let $E' \subset E$ be the differential graded $R$-subalgebra
with
$$
(E')^i = \left\{
\begin{matrix}
E^i & \text{if }i < 0 \\
\Ker(E^0 \to E^1) & \text{if }i = 0 \\
0 & \text{if }i > 0
\end{matrix}
\right.
$$
Then there are quasi-isomorphisms of differential graded
algebras $(A, \text{d}) \leftarrow (E', \text{d}) \rightarrow (E, \text{d})$.
Thus we obtain equivalences
$$
D(A, \text{d}) \leftarrow D(E', \text{d}) \rightarrow D(E, \text{d})
\rightarrow D(B, \text{d})
$$
by Lemma \ref{lemma-qis-equivalence}.
\end{proof}
\end{document}
